\documentclass[../constitution.tex]{subfiles}
\begin{document}

\hypertarget{part-5-committee}{%
  \part{COMMITTEE}\label{part-5-committee}}

\hypertarget{division-1-powers-of-committee}{%
  \section{Powers of Committee}\label{division-1-powers-of-committee}}

\hypertarget{committee}{%
  \subsection{Committee}\label{committee}}

\begin{enumerate}

  \item The committee members are the persons who, as the management committee of the Association, have the power to manage the affairs of the Association.
  \item Subject to the Act, these rules, the by-laws (if any) and any resolution passed at a general meeting, the committee has power to do all things necessary or convenient to be done for the proper management of the affairs of the Association.
  \item The committee must take all reasonable steps to ensure that the Association complies with the Act, these rules and the by-laws (if any).
\end{enumerate}

\hypertarget{division-2-composition-of-committee-and-duties-of-members}{%
  \section{Composition of committee and duties of members}\label{division-2-composition-of-committee-and-duties-of-members}}

\hypertarget{committee-members}{%
  \subsection{Committee members}\label{committee-members}}

\begin{enumerate}

  \item \label{number-of-committee-members} The committee must consist of at least 5 but not more than 10 committee members.
  \item \label{office-holders} The following are the office holders of the Association ---

        \begin{enumerate}

          \item the chairperson;
          \item the deputy chairperson;
          \item the secretary;
          \item the treasurer.
        \end{enumerate}
  \item A person may be a committee member if the person is --- \label{committee-eligibility}

        \begin{enumerate}

          \item an individual who has reached 18 years of age; and
          \item a ordinary member.
        \end{enumerate}
  \item A person must not hold 2 or more of the offices mentioned in subrule \ref{office-holders} at the same time.
\end{enumerate}

\hypertarget{chairperson}{%
  \subsection{Chairperson}\label{chairperson}}

\begin{enumerate}

  \item It is the duty of the chairperson to consult with the secretary regarding the business to be conducted at each committee meeting and general meeting.
  \item The chairperson has the powers and duties relating to convening and presiding at committee meetings and presiding at general meetings provided for in these rules.
\end{enumerate}

\hypertarget{secretary}{%
  \subsection{Secretary}\label{secretary}}

The Secretary has the following duties ---

\begin{enumerate}
  \item dealing with the Association's correspondence;
  \item consulting with the chairperson regarding the business to be conducted at each committee meeting and general meeting;
  \item preparing the notices required for meetings and for the business to be conducted at meetings;
  \item unless another member is authorised by the committee to do so, maintaining on behalf of the Association the register of members, and recording in the register any changes in the membership, as required under section 53(1) of the Act;
  \item maintaining on behalf of the Association an up-to-date copy of these rules, as required under section 35(1) of the Act;
  \item unless another member is authorised by the committee to do so, maintaining on behalf of the Association a record of committee members and other persons authorised to act on behalf of the Association, as required under section 58(2) of the Act;
  \item ensuring the safe custody of the books of the Association, other than the financial records, financial statements and financial reports, as applicable to the Association;
  \item maintaining full and accurate minutes of committee meetings and general meetings;
  \item carrying out any other duty given to the secretary under these rules or by the committee.
\end{enumerate}

\hypertarget{treasurer}{%
  \subsection{Treasurer}\label{treasurer}}

The treasurer has the following duties ---

\begin{enumerate}
  \item ensuring that any amounts payable to the Association are collected and issuing receipts for those amounts in the Association's name;
  \item ensuring that any amounts paid to the Association are credited to the appropriate account of the Association, as directed by the committee;
  \item ensuring that any payments to be made by the Association that have been authorised by the committee or at a general meeting are made on time;
  \item ensuring that the Association complies with the relevant requirements of Part 5 of the Act;
  \item ensuring the safe custody of the Association's financial records, financial statements and financial reports, as applicable to the Association;
  \item if the Association is a tier 1 association, coordinating the preparation of the Association's financial statements before their submission to the Association's annual general meeting;
  \item if the Association is a tier 2 association or tier 3 association, coordinating the preparation of the Association's financial report before its submission to the Association's annual general meeting;
  \item h providing any assistance required by an auditor or reviewer conducting an audit or review of the Association's financial statements or financial report under Part 5 Division 5 of the Act;

  \item carrying out any other duty given to the treasurer under these rules or by the committee.
\end{enumerate}

\hypertarget{division-3-election-of-committee-members-and-tenure-of-office}{%
  \section{Election of committee members and tenure of office}\label{division-3-election-of-committee-members-and-tenure-of-office}}

\hypertarget{how-members-become-committee-members}{%
  \subsection{How members become Committee members}\label{how-members-become-committee-members}}

A member becomes a committee member if the member ---

\begin{enumerate}

  \item is elected to the committee at a general meeting; or
  \item is appointed to the committee by the committee to fill a casual vacancy under rule \ref{filling-casual-vacancies}.
\end{enumerate}

\hypertarget{nomination-of-committee-members}{%
  \subsection{Nomination of committee members}\label{nomination-of-committee-members}}

\begin{enumerate}

  \item At least 21 days before an annual general meeting, the secretary must send written notice to all the members ---

        \begin{enumerate}

          \item calling for nominations for election to the committee; and
          \item stating the date by which nominations must be received by the secretary to comply with subrule \ref{notice-of-consideration-for-election-committee} and \ref{nomination-must-be-seconded}.
        \end{enumerate}
  \item A member who wishes to be considered for election to the committee at the annual general meeting must nominate for election by notifying the secretary at least 7 days before the annual general meeting. \label{notice-of-consideration-for-election-committee}
  \item Each nomination must be supported by a second member. The second member must notify the secretary of their support for the nomination at least 7 days before the annual general meeting. \label{nomination-must-be-seconded}
  \item Notices to the secretary for the purposes of rules 33.2 and 33.3 must be either ---

        \begin{enumerate}
          \item A written notice; or
          \item A direct message.
        \end{enumerate}

  \item (deleted.)
  \item A member whose nomination does not comply with this rule is not eligible for election to the committee unless the member is nominated under rule \ref{further-nominations-for-unfilled}.
\end{enumerate}

\hypertarget{election-of-committee-members}{%
  \subsection{Election of committee members}\label{election-of-committee-members}}

\begin{enumerate}

  \item No confidence in a nominee \label{no-confidence-in-committee-nominee}

        \begin{enumerate}
          \item At the annual general meeting, before the number of committee positions is decided under subrule \ref{decide-number-of-committee-members}, any ordinary member present at the meeting may move a motion of no confidence in a nominee for committee membership. The motion must be seconded by another ordinary member present. \label{raise-no-confidence-motion}
          \item If a motion is moved under subrule \ref{no-confidence-in-committee-nominee} ---
                \begin{enumerate}
                  \item If the nominee is present at the meeting, the chairperson must give the nominee up to 5 minutes to make oral representations in response to the motion. Those representations must not be defamatory.
                  \item If the nominee is not present at the meeting, the motion proceeds to a vote without any right of reply.
                  \item The meeting must vote on the motion by secret ballot.
                  \item The motion is carried if a majority of the ordinary members present and voting vote in favour of it.
                \end{enumerate}
          \item If a motion of no confidence in a nominee is carried, the nominee is removed from the list of nominees and does not stand for election at that meeting. \label{no-confidence-reduced-list-of-nominees}
          \item Subrules \ref{decide-number-of-committee-members}, \ref{nominating-committee-less-than-elected}, and \ref{nominating-committee-more-than-elected} apply after any motions under this rule have been decided, using the list of nominees as reduced by subrule \ref{no-confidence-reduced-list-of-nominees}.
          \item If a position for which a nominee was removed under this rule remains unfilled after subrules \ref{nominating-committee-less-than-elected} and \ref{nominating-committee-more-than-elected} are applied, that position is taken to be a casual vacancy for the purposes of rule \ref{filling-casual-vacancies}.
          \item The chairperson may decline to accept a motion under subrule \ref{raise-no-confidence-motion} if the chairperson considers it to be frivolous, vexatious, or dilatory. A member may move that the chairperson's decision under this subrule be overruled; that motion must be put to the meeting immediately and is carried by a majority of the ordinary members present and voting.
          \item A motion of no confidence must not be moved against the same nominee more than once at the same meeting.
        \end{enumerate}

  \item At the annual general meeting, the Association must decide by resolution the number of committee members to hold office for the next year. \label{decide-number-of-committee-members}

        \note {
          Refer to subrule \ref{number-of-committee-members} regarding the number of commitee members.
        }

  \item If the number of members nominating for committee membership is not greater than the number to be elected, the chairperson of the meeting --- \label{nominating-committee-less-than-elected}

        \begin{enumerate}

          \item must declare each of those members to be elected as commitee members; and \label{chairperson-declare-elected}
          \item may call for further nominations from the ordinary members at the meeting to fill any positions remaining unfilled after the elections under rule \ref{chairperson-declare-elected}. \label{further-nominations-for-unfilled}
        \end{enumerate}
  \item If --- \label{nominating-committee-more-than-elected}

        \begin{enumerate}

          \item the number of members nominating for committee membership is greater than the number to be elected; or
          \item the number of members nominating under rule \ref{further-nominations-for-unfilled} is greater than the number of positions remaining unfilled,
        \end{enumerate}

        the members at the meeting must vote by secret ballot to decide the members who are to be elected as committee members.

  \item A member who has nominated for committee membership may vote in accordance with that nomination.
\end{enumerate}

\hypertarget{election-of-office-holders}{%
  \subsection{Election of office holders}\label{election-of-office-holders}}

\begin{enumerate}

  \item Subject to subrule \ref{first-committee-meeting}, the newly elected committee members must hold the \textbf{first committee meeting} as soon as practicable after the annual general meeting.

  \item The chairperson of the first committee meeting must be ---

        \begin{enumerate}
          \item The first committee member to volunteer; or
          \item If no committee member volunteers - determined at random.
        \end{enumerate}

  \item The ordinary business of the first committee meeting is as follows ---

        \begin{enumerate}
          \item To elect the chairperson of the Association; then
          \item To elect the other office holders of the Association; then
          \item Any other business which the by-laws require to be addressed at the first committee meeting.
        \end{enumerate}




  \item A separate election must be held for each position of office holder of the Association, as follows ---

        \begin{enumerate}
          \item The chairperson of the meeting must call for nominations from the committee members at the meeting.
          \item If there is no nomination for a position, the chairperson of the meeting must appoint a committee member to the position.
          \item If only one committee member has nominated for a position, the chairperson of the meeting must declare the committee member elected to the position.
          \item If more than one committee member has nominated for a position, the committee members at the meeting must vote by secret ballot to decide who is to be elected to the position.
          \item Each committee member present at the meeting may vote for one committee member who has nominated for the position.
          \item A committee member who has nominated for the position may vote for themselves.
          \item On the committee member's election, the new chairperson of the Association may take over as the chairperson of the meeting.
        \end{enumerate}

\end{enumerate}

\hypertarget{term-of-office}{%
  \subsection{Term of office}\label{term-of-office}}

\begin{enumerate}

  \item The term of office of a committee member begins when the member ---

        \begin{enumerate}

          \item is elected at an annual general meeting or under rule \ref{elect-committee-at-general-meeting}; or
          \item is appointed to fill a casual vacancy under rule \ref{filling-casual-vacancies}.
        \end{enumerate}
  \item Subject to rule \ref{when-membership-of-committee-ceases}, a committee member holds office until the positions on the committee are declared vacant at the next annual general meeting.
  \item A committee member may be re-elected.
\end{enumerate}

\hypertarget{resignation-and-removal-from-office}{%
  \subsection{Resignation and removal from office}\label{resignation-and-removal-from-office}}

\begin{enumerate}

  \item A committee member may resign from the committee by written notice given to the secretary or, if the resigning member is the secretary, given to the chairperson.
  \item The resignation takes effect ---

        \begin{enumerate}

          \item when the notice is received by the secretary or chairperson; or
          \item if a later time is stated in the notice, at the later time.
        \end{enumerate}
  \item At a general meeting, the Association may by resolution ---

        \begin{enumerate}

          \item remove a committee member from office; and \label{remove-committee-member-resolution}
          \item elect a member who is eligible under rule \ref{committee-eligibility} to fill the vacant position. \label{elect-committee-at-general-meeting}
        \end{enumerate}
  \item A committee member who is the subject of a proposed resolution under subrule \ref{remove-committee-member-resolution} may make written representations (of a reasonable length) to the secretary or chairperson and may ask that the representations be provided to the members.
  \item The secretary or chairperson may give a copy of the representations to each member or, if they are not so given, the committee member may require them to be read out at the general meeting at which the resolution is to be considered.
\end{enumerate}

\hypertarget{when-membership-of-committee-ceases}{%
  \subsection{When membership of committee ceases}\label{when-membership-of-committee-ceases}}

\begin{enumerate}

  \item A person ceases to be a committee member if the person ---

        \begin{enumerate}

          \item dies or otherwise ceases to be a member; or
          \item resigns from the committee or is removed from office under rule \ref{resignation-and-removal-from-office}; or
          \item becomes ineligible to accept an appointment or act as a committee member under section 39 of the Act;
          \item becomes permanently unable to act as a committee member because of a mental or physical disability; or
          \item fails to attend 3 consecutive Committee meetings, of which the person has been given notice, without having notified the committee at least 24 hours beforehand that the person will be unable to attend.
        \end{enumerate}
\end{enumerate}

\note[Note for this rule]{
  Section 41 of the Act imposes requirements, arising when a person ceases to be a member of the management committee of an incorporated association, that relate to returning documents and records. }

\hypertarget{filling-casual-vacancies}{%
  \subsection{Filling casual vacancies}\label{filling-casual-vacancies}}

\begin{enumerate}

  \item The committee may appoint a member who is eligible under rule \ref{committee-eligibility} to fill a position on the committee that ---

        \begin{enumerate}

          \item has become vacant under rule \ref{when-membership-of-committee-ceases}; or
          \item was not filled by election at the most recent annual general meeting or under rule \ref{elect-committee-at-general-meeting}.
        \end{enumerate}
  \item If the position of secretary becomes vacant, the committee must appoint a member who is eligible under rule \ref{committee-eligibility} to fill the position within 14 days after the vacancy arises.
  \item Subject to the requirement for a quorum under rule \ref{quorum-for-committee-meetings}, the committee may continue to act despite any vacancy in its membership.
  \item If there are fewer committee members than required for a quorum under rule \ref{quorum-for-committee-meetings}, the committee may act only for the purpose of --- \label{committee-doesnt-meet-quorum}

        \begin{enumerate}
          \item appointing committee members under this rule; or
          \item convening a general meeting.
        \end{enumerate}
\end{enumerate}

\hypertarget{validity-of-acts}{%
  \subsection{Validity of acts}\label{validity-of-acts}}

The acts of a committee or subcommittee, or of a committee member or member of a subcommittee, are valid despite any defect that may afterwards be discovered in the election, appointment or qualification of a committee member or member of a subcommittee.

\hypertarget{payments-to-committee-members}{%
  \subsection{Payments to Committee Members}\label{payments-to-committee-members}}

\begin{enumerate}

  \item In this rule ---

        \textbf{committee member} includes a member of a subcommittee;

        \textbf{committee meeting} includes a meeting of a subcommittee.

  \item A committee member is entitled to be paid out of the funds of the Association for any out-of-pocket expenses for travel and accommodation incurred in connection with the Association's business, provided that ---

        \begin{enumerate}

          \item the expenses do \textbf{not} relate to attending a committee meeting; and
          \item the expenses do \textbf{not} relate to attending a general meeting; and
          \item The reason for the travel, and an estimate of the expenses, is approved by the committee before the expenses are incurred.
        \end{enumerate}
\end{enumerate}

\note[Note]{
  All committee members are also members of the Association. Therefore, committee members can also be paid out of the funds of the Association under subrule \ref{reasons-why-members-can-be-paid} regarding payments to members of the Association.}



\hypertarget{division-4-committee-meetings}{%
  \section{Committee meetings}\label{division-4-committee-meetings}}

\hypertarget{committee-meetings}{%
  \subsection{Committee meetings}\label{committee-meetings}}

\begin{enumerate}

  \item The committee will meet at least once every two months.
  \item \label{first-committee-meeting} The date, time and place of the first committee meeting must be determined by the committee members as soon as practicable after the annual general meeting at which the committee members are elected.

        \note{Refer to rule \ref{election-of-office-holders} for the business that must be considered at the first committee meeting.}

  \item Special committee meetings may be convened by the chairperson or any 3 committee members.
\end{enumerate}

\hypertarget{notice-of-committee-meetings}{%
  \subsection{Notice of committee meetings}\label{notice-of-committee-meetings}}

\begin{enumerate}

  \item Notice of each committee meeting must be given to each committee member at least 48 hours before the time of the meeting.
  \item The notice must state the date, time and place of the meeting (and, if the meeting is to be held in two or more places, the technology that will be used to facilitate this), and must describe the general nature of the business to be conducted at the meeting.
  \item Unless subrule \ref{unanimous-vote-urgent-business} applies, the only business that may be conducted at the meeting is the business described in the notice.
  \item Urgent business that has not been described in the notice may be conducted at the meeting if the committee members at the meeting unanimously agree to treat that business as urgent. \label{unanimous-vote-urgent-business}
\end{enumerate}

\hypertarget{procedure-and-order-of-business}{%
  \subsection{Procedure and order of business}\label{procedure-and-order-of-business}}

\begin{enumerate}

  \item The chairperson or, in the chairperson's absence, the deputy-chairperson must preside as chairperson of each committee meeting.
  \item If the chairperson and deputy chairperson are absent or are unwilling to act as chairperson of a meeting, the committee members at the meeting must choose one of them to act as chairperson of the meeting.
  \item The procedure to be followed at a committee meeting must be determined from time to time by the committee.
  \item The order of business at a committee meeting may be determined by the committee members at the meeting.
  \item A member or other person who is not a committee member may attend a committee meeting if invited to do so by the committee. \label{meeting-invitation-for-non-committee}
  \item A person invited under subrule \ref{meeting-invitation-for-non-committee} to attend a committee meeting ---

        \begin{enumerate}

          \item has no right to any agenda, minutes or other document circulated at the meeting; and
          \item must not comment about any matter discussed at the meeting unless invited by the committee to do so; and
          \item cannot vote on any matter that is to be decided at the meeting.
        \end{enumerate}
\end{enumerate}

\note{
  \textbf{Act Requirements - Material Personal Interests of Committee Members}

  \begin{itemize}
    \item Under section 42 of the Act a member of the committee who has a material personal interest in a matter being considered at a committee meeting must:
          \begin{itemize}
            \item as soon as he or she becomes aware of that interest, disclose the nature and extent of his or her interest to the Committee;
            \item disclose the nature and extent of the interest at the next general meeting of the association
          \end{itemize}
    \item Under section 42(3) of the Act this rule does not apply in respect of a material personal interest
          \begin{itemize}
            \item that exists only because the member ---
                  \begin{itemize}
                    \item is an employee of the incorporated association; or
                    \item is a member of a class of persons for whose benefit the association is established; or
                  \end{itemize}
            \item that the member has in common with all, or a substantial proportion of, the members of the Association.
          \end{itemize}
    \item Under section 43 of the Act a member of the management committee who has a material personal interest in a matter being considered at a meeting of the management committee must not be present while the matter is being considered at the meeting or vote on the matter.
  \end{itemize}

  Under section 42(6) of the Act the association must record every disclosure made by a committee member of a material personal interest in the minutes of the committee meeting at which the disclosure is made.

}

\hypertarget{use-of-technology-to-be-present-at-committee-meetings}{%
  \subsection{Use of technology to be present at committee meetings}\label{use-of-technology-to-be-present-at-committee-meetings}}

If a committee meeting is held in two or more places linked together by any technology ---

\begin{enumerate}

  \item a member present at one of the places is taken to be present at the meeting unless and until that member states to the chair of the meeting that they are discontinuing participation in the meeting; and

  \item the chair of that meeting may determine at which place the meeting will be taken to have been held.

\end{enumerate}



\hypertarget{quorum-for-committee-meetings}{%
  \subsection{Quorum for committee meetings}\label{quorum-for-committee-meetings}}

\begin{enumerate}

  \item Subject to subrule \ref{committee-doesnt-meet-quorum} no business is to be conducted at a committee meeting unless a quorum is present.
  \item If a quorum is not present within 30 minutes after the notified commencement time of a committee meeting --- \label{no-quorum-after-30min}

        \begin{enumerate}
          \item in the case of a special meeting --- the meeting lapses; or
          \item otherwise, the meeting is adjourned to the same time, day and place in the following week. \label{meeting-adjourned-due-to-no-quorum}
        \end{enumerate}
  \item If ---

        \begin{enumerate}

          \item a quorum is not present within 30 minutes after the commencement time of a committee meeting held under subrule \ref{meeting-adjourned-due-to-no-quorum}; and
          \item at least 2 committee members are present at the meeting,
        \end{enumerate}

        those members present are taken to constitute a quorum.

\end{enumerate}

\hypertarget{voting-at-committee-meetings}{%
  \subsection{Voting at committee meetings}\label{voting-at-committee-meetings}}

\begin{enumerate}

  \item Each committee member present at a committee meeting has one vote on any question arising at the meeting.
  \item A motion is carried if a majority of the committee members present at the committee meeting vote in favour of the motion.
  \item If the votes are divided equally on a question, the chairperson of the meeting has a second or casting vote.
  \item A vote may take place by the committee members present indicating their agreement or disagreement or by a show of hands, unless the committee decides that a secret ballot is needed to determine a particular question.
  \item If a secret ballot is needed, the chairperson of the meeting must decide how the ballot is to be conducted.
\end{enumerate}

\hypertarget{minutes-of-committee-meetings}{%
  \subsection{Minutes of committee meetings}\label{minutes-of-committee-meetings}}

\begin{enumerate}

  \item The committee must ensure that minutes are taken and kept of each committee meeting.
  \item The minutes must record the following ---

        \begin{enumerate}

          \item the names of the committee members present at the meeting;
          \item the name of any person attending the meeting under rule \ref{meeting-invitation-for-non-committee};
          \item the business considered at the meeting;
          \item any motion on which a vote is taken at the meeting and the result of the vote.
        \end{enumerate}
  \item The minutes of a committee meeting must be entered in the Association's minute book within 7 days after the meeting is held.
  \item The committee must ensure that the minutes of a committee meeting are reviewed and approved as correct by a motion at the next committee meeting.
  \item When the minutes of a committee meeting have been accepted as correct they are, until the contrary is proved, evidence that ---

        \begin{enumerate}

          \item the meeting to which the minutes relate was duly convened and held; and
          \item the matters recorded as having taken place at the meeting took place as recorded; and
          \item any appointment purportedly made at the meeting was validly made.
        \end{enumerate}
\end{enumerate}

\note[Note for this rule]{

  Section 42(6) of the Act requires details relating to the disclosure of a committee member's material personal interest in a matter being considered at a committee meeting to be recorded in the minutes of the meeting.

}

\hypertarget{division-5-subcommittees-and-subsidiary-offices}{%
  \section{Subcommittees and subsidiary offices}\label{division-5-subcommittees-and-subsidiary-offices}}

\hypertarget{subcommittees-and-subsidiary-offices}{%
  \subsection{Subcommittees and subsidiary offices}\label{subcommittees-and-subsidiary-offices}}

\begin{enumerate}

  \item To help the committee in the conduct of the Association's business, the committee may, in writing, do either or both of the following ---

        \begin{enumerate}

          \item appoint one or more subcommittees; \label{appoint-subcommittees}
          \item create one or more subsidiary offices and appoint people to those offices.
        \end{enumerate}
  \item A subcommittee may consist of the number of people, whether or not members, that the committee considers appropriate.
  \item A person may be appointed to a subsidiary office whether or not the person is a member.
  \item Subject to any directions given by the committee ---

        \begin{enumerate}

          \item a subcommittee may meet and conduct business as it considers appropriate; and
          \item the holder of a subsidiary office may carry out the functions given to the holder as the holder considers appropriate.
        \end{enumerate}
\end{enumerate}

\hypertarget{delegation-to-subcommittees-and-holders-of-subsidiary-offices}{%
  \subsection{Delegation to subcommittees and holders of subsidiary offices}\label{delegation-to-subcommittees-and-holders-of-subsidiary-offices}}

\begin{enumerate}

  \item In this rule ---

        \textbf{non-delegable duty} means a duty imposed on the committee by the Act or another written law.

  \item The committee may, in writing, delegate to a subcommittee or the holder of a subsidiary office the exercise of any power or the performance of any duty of the committee other than ---

        \begin{enumerate}

          \item the power to delegate; and
          \item a non-delegable duty.
        \end{enumerate}
  \item A power or duty, the exercise or performance of which has been delegated to a subcommittee or the holder of a subsidiary office under this rule, may be exercised or performed by the subcommittee or holder in accordance with the terms of the delegation.
  \item The delegation may be made subject to any conditions, qualifications, limitations or exceptions that the committee specifies in the document by which the delegation is made.
  \item The delegation does not prevent the committee from exercising or performing at any time the power or duty delegated.
  \item Any act or thing done by a subcommittee or by the holder of a subsidiary office, under the delegation has the same force and effect as if it had been done by the committee.
  \item The committee may, in writing, amend or revoke the delegation.
\end{enumerate}

\end{document}